\begin{exercisebox}[label={box:LOalphas}]{LO $\alpha_s$ evolution 1}
    Proof \cref{eq:LOalphas} can be written as
    \begin{equation}
        \alpha_s(\mu_R^2) = \frac 1 {\beta_0 \ln(\mu_R^2/\LQCD^2)}
    \end{equation}
    and given an explicit expression for $\LQCD$.
\end{exercisebox}

Before we go to NLO let's use the more practical notation:
\begin{exercisebox}{LO $\alpha_s$ evolution 2}
    Show
    \begin{equation}
        \alpha_s^\text{LO}(t_R) = \frac{\alpha_s^\text{LO}(0)}{1 + \alpha_s^\text{LO}(0) \beta_0 t_R}
    \end{equation}
    (i.e.\ \cref{eq:LOalphas} recast in terms of $t_R = \ln(\mu_R^2/\mu_{R,0}^2)$) is a solution to
    \begin{equation}
        \dv {\alpha_s^\text{LO}(t_R)}{t_R} = -\beta_0 \left(\alpha_s^\text{LO}(t_R)\right)^2
    \end{equation}
    (i.e.\ \cref{eq:beta} at LO recast in terms of $t_R$).
\end{exercisebox}

\begin{exercisebox}[label={box:NLOalphas}]{NLO $\alpha_s$ evolution 1}
    Show
    \begin{equation}
        \alpha_s^\text{NLO}(t_R) = \alpha_s^\text{LO}(t_R) - \frac{\beta_1}{\beta_0} \left(\alpha_s^\text{LO}(t_R)\right)^2 \ln(\alpha_s^\text{NLO}(0) / \alpha_s^\text{LO}(t_R)) \label{eq:NLOalphast}
    \end{equation}
    is a solution to
    \begin{equation}
        \dv {\alpha_s^\text{NLO}(t_R)}{t_R} = -\beta_0 \left(\alpha_s^\text{NLO}(t_R)\right)^2 -\beta_1 \left(\alpha_s^\text{NLO}(t_R)\right)^3 \label{eq:NLObetat}\,
    \end{equation}
    (i.e.\ \cref{eq:NLObeta} recast in terms of $t_R$) at NLO accuracy, i.e.\ the full solution (which is not known analytically) and $\alpha_s^\text{NLO}(t_R)$ differ by higher order terms.
\end{exercisebox}
\noindent Make two steps:
\begin{enumerate}
    \item compute the left-hand-side of \cref{eq:NLObetat} using \cref{eq:NLOalphast}
    \item compute the right-hand-side of \cref{eq:NLObetat} using \cref{eq:NLOalphast}
\end{enumerate}
Show the two sides differ by $\order{\qty(\alpha_s(t_R))^4}$, i.e.\ they agree up to the required accuracy.

\begin{exercisebox}[label={box:NLOalphasLL}]{NLO $\alpha_s$ evolution 2}
    Show that \cref{eq:NLObetat} performs a next-to-leading logarithmic (NLL) resummation.
\end{exercisebox}
\noindent The leading logarithmic (LL) resummation is visible in \cref{eq:LOalphasresum}, where for each logarithm there is the adjacent coupling, in other words: for each coupling $\alpha_s(\mu_{R,0})^k$, there is one $\ln^k(\mu_R^2/\mu_{R,0}^2)$ or, equivalently, one $t_R^k$.
At NLL we must in addition resum also terms $\alpha_s(\mu_{R,0})^k t_R^{k-1}$, i.e.\ where the logarithm has one power less.
Note that \cref{eq:LOalphasresum} has an overall factor of $\alpha_s(\mu_{R,0})$, which we neglect here since the resummation happens for the solution operator $T$, which we can define in an analogous way to the EKO, \cref{eq:EKO},
\begin{equation}
    \alpha_s(\mu_R^2) = T(\mu_R^2 \leftarrow \mu_{R,0}^2) \alpha_s(\mu_{R,0}^2)\,.
\end{equation}

Make several steps:
\begin{enumerate}
    \item Expand $\alpha_s^\text{LO}(t_R)$ in $t_R$ up to and including $\order{t_R^2}$ using \cref{eq:LOalphasresum}
    \item Insert this expansion into \cref{eq:NLOalphast} and expand $\alpha_s^\text{NLO}(t_R)$ to the same accuracy. Recall $\ln(1 + w \epsilon) = w \epsilon - \frac 1 2 w^2 \epsilon^2 + \order{\epsilon^3}$.
    \item Independently we make a Taylor expansion of $\alpha_s^\text{NLO}(t_R)$ up to and including $\order{t_R^2}$.
        The first derivative is \cref{eq:NLObetat}.
        Compute the second derivative using \cref{eq:NLObetat}.
    \item Check that the two expression agree at LL, $\alpha_s(\mu_{R,0})^k t_R^{k}$ for $k=0,1,2$
    \item Check that the two expression agree at NLL, $\alpha_s(\mu_{R,0})^k t_R^{k-1}$ for $k=1,2,3$
    \item Check that the two expression start to disagree at next-to-NLL (NNLL), $\alpha_s(\mu_{R,0})^k t_R^{k-2}$ for $k=2,3,4$
\end{enumerate}
